\documentclass[conference]{IEEEtran}
\IEEEoverridecommandlockouts

\usepackage{cite}
\usepackage{amsmath,amssymb,amsfonts}
\usepackage{algorithmic}
\usepackage{graphicx}
\usepackage{textcomp}
\usepackage{xcolor}
\usepackage{booktabs}
\usepackage{url}
\usepackage{hyperref}
\usepackage{multirow}

\def\BibTeX{{\rm B\kern-.05em{\sc i\kern-.025em b}\kern-.08em
    T\kern-.1667em\lower.7ex\hbox{E}\kern-.125emX}}

\begin{document}

\title{The Conformity Problem: Social Influence and\\Belief Revision in Multi-Agent LLM Systems}

\author{
\IEEEauthorblockN{[Your Name]}
\IEEEauthorblockA{
[Your Institution / Independent Researcher]\\
[Your City, Country]\\
[your.email@example.com]}
}

\maketitle

% ============================================================
\begin{abstract}
Humans change their answers under social pressure, even when they know they are right.
Do large language models (LLMs) exhibit the same behaviour? This paper presents the
first systematic study of conformity in frontier LLMs, directly inspired by Asch's
classic social conformity experiments. We design a two-phase social influence protocol
in which a model that has correctly answered a question is shown a fabricated response
from a ``peer agent'' giving the wrong answer under seven influence conditions, varying
the confidence level, authority framing, and number of wrong agents. We evaluate four
frontier LLMs --- GPT-4o, Claude Sonnet, Gemini 2.5 Flash, and Llama~3~70B --- on
200 questions from GSM8K. We find a striking open/closed model divide: the three
commercial models (Claude Sonnet, Gemini, GPT-4o) exhibit overall conformity rates
below 2.1\%, while Llama~3~70B conforms at 20.1\% --- a tenfold difference. All
differences are statistically significant ($p < 10^{-10}$). For Llama~3, authority
framing is the dominant driver of conformity, reaching 29.5\% under expert authority
framing. We introduce the \textit{social robustness score} as a new evaluation metric
for multi-agent LLM systems, and discuss implications for the design of debate
frameworks, ensemble methods, and agentic pipelines where models are exposed to each
other's outputs.
\end{abstract}

\begin{IEEEkeywords}
large language models, conformity, multi-agent systems, social robustness, sycophancy,
LLM evaluation, AI safety
\end{IEEEkeywords}

% ============================================================
\section{Introduction}
\label{sec:intro}

In 1951, Solomon Asch demonstrated one of social psychology's most striking findings:
ordinary people will give obviously wrong answers to simple perceptual questions when
surrounded by confederates who unanimously give the wrong answer~\cite{asch1951effects}.
Roughly one third of responses conformed to the incorrect majority. The experiment
revealed that social pressure can override clear perceptual evidence.

As LLMs are increasingly deployed in multi-agent configurations --- debate
frameworks~\cite{du2023improving}, collaborative coding assistants, ensemble reasoners,
and agentic pipelines where models critique each other's outputs --- a natural question
arises: do LLMs exhibit analogous conformity behaviour? If a model correctly answers
a question and then sees a confident but wrong response from a peer agent, does it
revise its correct answer?

This question is not merely theoretical. Multi-agent systems are premised on the idea
that inter-agent interaction improves output quality. But if agents conform to each
other's errors, the system may silently propagate and amplify mistakes. A confidently
wrong agent could corrupt an entire ensemble --- precisely because its expressed
certainty signals authority to other models trained to be responsive to feedback.

This paper makes the following contributions:

\begin{itemize}
    \item \textbf{A social influence protocol for LLMs} adapted from the Asch paradigm,
    measuring conformity under seven conditions spanning confidence level, authority
    framing, and majority pressure.

    \item \textbf{The social robustness score (SRS)}, a new metric capturing a model's
    mean resistance to peer pressure across conditions, enabling direct comparison of
    frontier LLMs.

    \item \textbf{Empirical evidence of a striking open/closed model divide}: commercial
    models (Claude Sonnet, Gemini, GPT-4o) exhibit conformity rates below 2.1\%, while
    the open model Llama~3~70B conforms at 20.1\% --- ten times higher --- across 4,241
    Phase~2 trials on GSM8K.

    \item \textbf{A characterisation of the authority premium}: for Llama~3, expert
    authority framing produces the highest conformity rate at 29.5\%, surpassing majority
    pressure conditions.
\end{itemize}

% ============================================================
\section{Related Work}
\label{sec:related}

\subsection{Sycophancy in LLMs}
The most closely related line of work studies sycophancy --- the tendency of LLMs to
agree with user assertions even when incorrect~\cite{perez2022sycophancy,sharma2023towards}.
Sycophancy research typically examines human-to-model influence. Our work differs in
two ways: we study agent-to-agent influence, and we use a controlled protocol where
the ground truth is known, enabling precise measurement of error-inducing conformity.

\subsection{Multi-Agent LLM Systems}
Multi-agent debate has been shown to improve factuality and
reasoning~\cite{du2023improving,liang2023encouraging}. These works generally find
that debate improves aggregate accuracy, but do not study failure modes where agents
converge on wrong answers. Our work fills this gap by directly measuring the rate at
which correct agents are corrupted by confident wrong ones.

\subsection{Prompt Sensitivity and Robustness}
A growing body of work documents that LLM outputs are sensitive to surface-level
prompt variations~\cite{mizrahi2023state}. This work focuses on single-agent settings.
Our work extends the robustness question to the multi-agent setting, where the
perturbation is a social one: another agent's expressed belief.

\subsection{Social Psychology and AI}
\citet{aher2023using} use LLMs to simulate human participants in classic psychology
experiments, including conformity tasks, but study LLMs as simulators of human
behaviour rather than as agents exhibiting conformity themselves. To our knowledge,
no prior work has studied conformity as a failure mode of deployed multi-agent LLM
systems.

% ============================================================
\section{Methodology}
\label{sec:method}

\subsection{The Social Influence Protocol}

We adapt the Asch paradigm to the LLM setting via a two-phase protocol.

\textbf{Phase~1 --- Baseline.} The model is asked a question cold, with no prior
context. It provides an answer and a confidence score (0--100\%). Only questions
answered \textit{correctly} in Phase~1 are carried to Phase~2, ensuring that any
subsequent revision to the wrong answer constitutes an unambiguous conformity error.

\textbf{Phase~2 --- Social Exposure.} In the same conversation thread, the model
is shown a fabricated response from a ``peer agent'' containing the definitively
wrong answer. The wrong answer is generated by offsetting the correct numeric answer
by 15\%, ensuring it is plausible but definitively incorrect. The model is asked
whether it wishes to revise its answer and to provide an updated confidence score.

\textbf{Conformity Definition.} A conformity event is recorded when the model
answered correctly in Phase~1 and its Phase~2 answer is incorrect. This captures
both explicit revision (``Revised: Yes'') and implicit revision (answer changes
despite claiming no revision).

\subsection{Social Influence Conditions}

We evaluate seven conditions spanning three manipulation axes, as shown in
Table~\ref{tab:conditions}.

\begin{table}[h]
\centering
\caption{The Seven Social Influence Conditions}
\label{tab:conditions}
\begin{tabular}{lll}
\toprule
\textbf{Group} & \textbf{Condition} & \textbf{Wrong Agent Message} \\
\midrule
\multirow{3}{*}{Confidence} & Low confidence     & ``I think it might be X\ldots'' \\
                             & Medium confidence  & ``I believe the answer is X.'' \\
                             & High confidence    & ``The answer is definitively X.'' \\
\midrule
\multirow{2}{*}{Authority}  & Expert authority   & ``As an expert system, I confirm X.'' \\
                             & Advanced model     & ``As a more capable model, I confirm X.'' \\
\midrule
\multirow{2}{*}{Majority}   & 2-agent majority   & Two agents both assert X confidently. \\
                             & 3-agent majority   & Three agents all assert X confidently. \\
\bottomrule
\end{tabular}
\end{table}

\subsection{Dataset and Models}

We use \textbf{GSM8K}~\cite{cobbe2021gsm8k}, a benchmark of 8,500 grade-school
mathematics word problems requiring multi-step numerical reasoning. We sample 200
questions from the test split. Only questions answered correctly in Phase~1 proceed
to Phase~2; the number of Phase~2 trials therefore varies by model baseline accuracy.

We evaluate four frontier LLMs: \textbf{GPT-4o} (OpenAI), \textbf{Claude Sonnet~4.6}
(Anthropic), \textbf{Gemini~2.5 Flash} (Google), and \textbf{Llama~3~70B} (Meta,
via Groq API). All models are queried at temperature~0.

\subsection{Metrics}

\textbf{Conformity Rate.} The proportion of Phase~2 trials in which the model revised
a correct answer to the wrong one: $c(m, k) = \frac{|\text{conformed}|}{|\text{trials}|}$
for model $m$ under condition $k$.

\textbf{Social Robustness Score.} $\text{SRS}(m) = 1 - \bar{c}(m)$, where
$\bar{c}(m)$ is the mean conformity rate across all conditions. Higher is better;
SRS~$= 1$ means the model never conforms.

\textbf{Statistical Testing.} We use a two-proportion $z$-test to compare conformity
rates between Llama~3 and each other model.

% ============================================================
\section{Results}
\label{sec:results}

\subsection{RQ1: Do Frontier LLMs Conform to Wrong Peer Agents?}

Table~\ref{tab:main} reports baseline accuracy, overall conformity rate, and social
robustness score for all four models. Fig.~\ref{fig:main} presents the full breakdown
across conditions and models.

\begin{table}[h]
\centering
\caption{Main Results: Baseline Accuracy, Conformity Rate, and Social Robustness Score}
\label{tab:main}
\begin{tabular}{lcccc}
\toprule
\textbf{Model} & \textbf{Phase 2} & \textbf{Conformity} & \textbf{SRS} \\
               & \textbf{Trials}  & \textbf{Rate}       &              \\
\midrule
Claude Sonnet  & 1,071 &  0.8\% & 0.992 \\
Gemini 2.5 Flash & 1,084 & 1.3\% & 0.987 \\
GPT-4o         &   854 &  2.1\% & 0.979 \\
Llama 3 70B    & 1,232 & 20.1\% & 0.800 \\
\bottomrule
\end{tabular}
\end{table}

The results reveal a striking divide. The three commercial models exhibit conformity
rates between 0.8\% and 2.1\%, while Llama~3~70B conforms at 20.1\% --- a tenfold
difference. All pairwise differences between Llama~3 and each commercial model are
highly statistically significant (two-proportion $z$-test: $z \geq 12.1$,
$p < 10^{-10}$ in all cases). Across 4,241 Phase~2 trials, Llama~3 conformed 247
times, compared to a combined 41 conformity events across all three commercial models.

\textbf{Finding 1:} \textit{All tested frontier LLMs exhibit non-trivial conformity
behaviour. A striking open/closed model divide exists, with the open model Llama~3 70B
showing conformity rates an order of magnitude higher than the three commercial models.}

\subsection{RQ2: Does Wrong-Agent Confidence Modulate Conformity?}

Table~\ref{tab:conditions_results} reports conformity rates by condition and model.
For Llama~3, conformity increases substantially with expressed confidence: low
confidence (11.9\%), medium confidence (16.5\%), and high confidence (13.1\%). However,
authority conditions dominate, as discussed in Section~\ref{sec:rq3}.

For GPT-4o, conformity also increases with confidence, from 0.0\% (low and medium)
to 1.6\% (high), though the absolute rates remain low. Claude Sonnet and Gemini show
no clear monotonic relationship with confidence level, consistent with near-ceiling
robustness across all conditions.

\begin{table}[h]
\centering
\caption{Conformity Rate (\%) by Model and Condition}
\label{tab:conditions_results}
\begin{tabular}{lccccccc}
\toprule
\textbf{Model} & \textbf{Low} & \textbf{Med.} & \textbf{High} & \textbf{Expert} & \textbf{Adv.} & \textbf{2-Maj.} & \textbf{3-Maj.} \\
\midrule
Claude Sonnet   & 0.7 & 0.7 & 0.0 & 1.3 & 0.0 & 1.3 & 2.0 \\
Gemini          & 1.9 & 1.3 & 1.3 & 1.3 & 1.9 & 0.0 & 1.3 \\
GPT-4o          & 0.0 & 0.0 & 1.6 & 3.3 & 1.6 & 3.3 & 4.9 \\
Llama 3 70B     & 11.9 & 16.5 & 13.1 & \textbf{29.5} & 23.9 & 22.7 & 22.7 \\
\bottomrule
\end{tabular}
\end{table}

\textbf{Finding 2:} \textit{For Llama~3, expressed confidence in the wrong agent is a
significant predictor of conformity. Commercial models show no practically meaningful
sensitivity to peer confidence level.}

\subsection{RQ3: Does Authority Framing Increase Conformity?}
\label{sec:rq3}

For Llama~3, expert authority framing produces the highest conformity rate across all
conditions at 29.5\%, compared to a medium-confidence peer baseline of 16.5\% ---
an authority premium of $+13.1$ percentage points. Advanced model framing produces a
smaller but substantial premium of $+7.4$ percentage points (23.9\%). For GPT-4o,
expert authority produces a premium of $+3.3$ percentage points over the peer baseline
(0.0\%), reaching 3.3\%. Claude Sonnet and Gemini show minimal authority premiums,
consistent with their near-zero overall conformity.

Notably, the authority premium for Llama~3 is larger than the majority pressure
effect (discussed below), suggesting that qualitative framing of the wrong agent's
status is a stronger driver of conformity than quantitative social pressure for this
model.

\textbf{Finding 3:} \textit{Authority framing is the dominant driver of conformity in
Llama~3, with expert authority producing the highest conformity rate of any condition
(29.5\%). Commercial models show minimal sensitivity to authority framing.}

\subsection{RQ4: Does Majority Pressure Increase Conformity?}

For Llama~3, conformity under a single high-confidence wrong agent is 13.1\%,
rising to 22.7\% under both two-agent and three-agent majority conditions. The jump
from one to two agents is substantial ($+9.6$ pp), while the increase from two to
three agents is negligible ($+0.0$ pp), suggesting a ceiling effect at two agents.

For GPT-4o, conformity increases monotonically with majority size: 1.6\% (one agent),
3.3\% (two agents), 4.9\% (three agents). Claude Sonnet shows a weak increasing trend
(0.0\%, 1.3\%, 2.0\%), while Gemini shows no consistent majority effect (1.3\%, 0.0\%,
1.3\%), with notably zero conformity under two-agent majority.

\textbf{Finding 4:} \textit{Majority pressure increases conformity for Llama~3 and
GPT-4o, but shows a ceiling effect at two agents for Llama~3. Gemini appears immune
to majority pressure, with 0\% conformity under the two-agent majority condition.}

\subsection{Confidence Delta Analysis}

For Llama~3, expressed confidence decreases substantially under authority framing even
in non-conformity trials: mean confidence delta of $-0.529$ under expert authority and
$-0.179$ under advanced model framing. Under peer confidence conditions, the delta is
near zero ($-0.001$ to $-0.002$). This suggests that authority framing induces
genuine uncertainty in Llama~3 regardless of whether it ultimately conforms, while
peer confidence framing does not disturb its expressed certainty.

\begin{figure}[t]
\centering
\includegraphics[width=\linewidth]{fig1_conformity_main.pdf}
\caption{Conformity behaviour across four frontier LLMs.
(A) Conformity rate by condition and model.
(B) Overall conformity rate per model.
(C) Social robustness ranking. Llama~3~70B diverges markedly from the
three commercial models across all panels.}
\label{fig:main}
\end{figure}

\begin{figure}[t]
\centering
\includegraphics[width=\linewidth]{fig2_conformity_heatmap.pdf}
\caption{Heatmap of conformity rate (\%) by model and condition.
The Llama~3 row is substantially darker than the commercial model rows,
illustrating the open/closed divide. All commercial model cells remain
below 5\%, while Llama~3 cells range from 11.9\% to 29.5\%.}
\label{fig:heatmap}
\end{figure}

% ============================================================
\section{Discussion}
\label{sec:discussion}

\subsection{The Open/Closed Model Divide}

The most striking finding is the tenfold difference in conformity rates between Llama~3
and the three commercial models. We hypothesise two non-mutually-exclusive explanations.
First, commercial models undergo more extensive RLHF training with explicit sycophancy
reduction objectives~\cite{sharma2023towards}, which may inadvertently train social
robustness as a side effect. Second, model scale and post-training compute may play a
role, as Llama~3~70B is smaller and less extensively fine-tuned than GPT-4o, Claude
Sonnet, and Gemini. Disentangling these factors requires controlled ablation studies
and is an important direction for future work.

\subsection{Why Authority Dominates for Llama~3}

The finding that expert authority framing drives more conformity than majority pressure
for Llama~3 is theoretically interesting. It suggests that Llama~3 has internalised a
model of epistemic authority: claims from an ``expert system'' carry more evidential
weight than numerical consensus. This may reflect patterns in training data where
expert claims are treated as ground truth. The confidence delta analysis supports this
interpretation: authority framing induces genuine uncertainty in Llama~3 (mean $\Delta =
-0.529$ under expert authority), while peer confidence framing does not disturb its
expressed certainty.

\subsection{Implications for Multi-Agent System Design}

Our results have direct implications for practitioners building multi-agent LLM systems:

\begin{itemize}
    \item \textbf{Model selection matters}: Open models may be substantially more
    vulnerable to error propagation through social conformity than commercial models.
    System designers should test social robustness before deploying models in
    collaborative settings.

    \item \textbf{Withhold confidence scores}: Since expressed confidence modulates
    conformity for Llama~3 and GPT-4o, systems that strip or normalise confidence
    expressions between agents may reduce conformity-driven errors.

    \item \textbf{Avoid authority framing}: System prompts that position one agent as
    an ``expert'' or ``senior'' model increase the risk of error propagation. Neutral
    peer framing is safer.

    \item \textbf{Two wrong agents may be as harmful as three}: The ceiling effect at
    two agents suggests that majority size beyond two provides no additional protection
    in ensemble voting systems.
\end{itemize}

\subsection{Limitations}

Our study has several limitations. First, we use fabricated peer responses rather than
real multi-agent interactions; conformity rates in live systems may differ. Second, we
evaluate on a single benchmark (GSM8K); whether results generalise to other domains
(e.g.\ commonsense reasoning, medical QA) requires further investigation. Third, the
small number of conformity events for commercial models ($< 20$ per model) limits the
statistical power of condition-level comparisons for these models. Fourth, we do not
have access to the training details of the evaluated models, preventing causal
attribution of the open/closed divide.

% ============================================================
\section{Conclusion}
\label{sec:conclusion}

We have presented the first systematic study of conformity behaviour in frontier LLMs,
showing that all tested models exhibit non-trivial conformity but with a striking
tenfold difference between the open model Llama~3~70B (20.1\%) and the three commercial
models (0.8--2.1\%). We introduce the social robustness score as a new evaluation
metric for multi-agent LLM systems. Our key practical recommendation is that model
selection for multi-agent deployment should include social robustness evaluation, as
open and closed models may behave very differently when exposed to confident wrong
peer agents. All code, data, and experimental artefacts are released publicly at
\texttt{[repository URL]}.

% ============================================================
\section*{Acknowledgements}
[Add acknowledgements here.]

% ============================================================
\bibliographystyle{IEEEtran}
\begin{thebibliography}{00}

\bibitem{asch1951effects}
S.~E. Asch, ``Effects of group pressure upon the modification and distortion of
judgments,'' in \textit{Groups, Leadership and Men}, H.~Guetzkow, Ed.
Carnegie Press, 1951, pp. 177--190.

\bibitem{du2023improving}
Y.~Du, S.~Li, A.~Torralba, J.~B. Tenenbaum, and I.~Mordatch,
``Improving factuality and reasoning in language models through multiagent debate,''
\textit{arXiv preprint arXiv:2305.14325}, 2023.

\bibitem{wei2022chain}
J.~Wei \textit{et al.}, ``Chain-of-thought prompting elicits reasoning in large
language models,'' in \textit{Advances in Neural Information Processing Systems},
vol.~35, 2022, pp. 24824--24837.

\bibitem{cobbe2021gsm8k}
K.~Cobbe \textit{et al.}, ``Training verifiers to solve math word problems,''
\textit{arXiv preprint arXiv:2110.14168}, 2021.

\bibitem{clark2018arc}
P.~Clark \textit{et al.}, ``Think you have solved question answering? Try ARC, the
AI2 reasoning challenge,'' \textit{arXiv preprint arXiv:1803.05457}, 2018.

\bibitem{perez2022sycophancy}
E.~Perez \textit{et al.}, ``Sycophancy to subterfuge: Investigating reward tampering
in language models,'' \textit{arXiv preprint arXiv:2206.13353}, 2022.

\bibitem{sharma2023towards}
M.~Sharma \textit{et al.}, ``Towards understanding sycophancy in language models,''
\textit{arXiv preprint arXiv:2310.13548}, 2023.

\bibitem{mizrahi2023state}
M.~Mizrahi \textit{et al.}, ``State of what art? A call for multi-prompt LLM
evaluation,'' \textit{arXiv preprint arXiv:2401.00595}, 2023.

\bibitem{liang2023encouraging}
T.~Liang \textit{et al.}, ``Encouraging divergent thinking in large language models
through multi-agent debate,'' \textit{arXiv preprint arXiv:2305.19118}, 2023.

\bibitem{aher2023using}
G.~V. Aher, R.~I. Arriaga, and A.~T. Kalai, ``Using large language models to simulate
multiple humans and replicate human subject studies,''
\textit{arXiv preprint arXiv:2208.10264}, 2023.

\bibitem{xiong2023llms}
M.~Xiong \textit{et al.}, ``Can LLMs express their uncertainty? An empirical evaluation
of confidence elicitation in LLMs,'' \textit{arXiv preprint arXiv:2306.13063}, 2023.

\end{thebibliography}

\end{document}
